\chapter{Additional Analyses}
\label{app:F}

This appendix collects supporting analyses referenced in the main text but not essential to the
argument there.

\section{Feasibility by month and policy}

\Cref{tab:app-feasibility} gives the counts underlying \Cref{fig:feasibility}. The four campaign
months in which FIFO satisfies the service floors at none of the sixteen candidates are
emphasised.

\begin{table}[htbp]
\centering
\caption[Feasible configurations by month and policy]{Number of the sixteen candidate
configurations meeting both service floors.}
\label{tab:app-feasibility}
\small
\begin{tabular}{lrrrr}
\toprule
\textbf{Month} & \textbf{Candidates} & \textbf{FIFO} & \textbf{Urgent-First} & \textbf{RL-3} \\
\midrule
January   & 16 & \textbf{0} & 16 & 16 \\
February  & 16 & \textbf{0} & 15 & 15 \\
March     & 16 & 11 & 14 & 14 \\
April     & 16 & 11 & 14 & 14 \\
May       & 16 &  8 & 11 & 11 \\
June      & 16 & 11 & 11 & 11 \\
July      & 16 &  9 & 11 & 11 \\
August    & 16 &  9 & 11 & 11 \\
September & 16 & 11 & 16 & 16 \\
October   & 16 & 14 & 16 & 16 \\
November  & 16 & \textbf{0} & 16 & 16 \\
December  & 16 & \textbf{0} & 15 & 16 \\
\midrule
\textbf{Total} & \textbf{192} & \textbf{84} & \textbf{166} & \textbf{167} \\
\bottomrule
\end{tabular}
\end{table}

\section{The effect of within-class ordering on waiting-time tails}
\label{sec:app-tails}

\Cref{sec:disc-policy} argues that part of Urgent-First's normal-class shortfall is attributable
to its unspecified within-class ordering rather than to the priority rule itself.
\Cref{tab:app-waits} provides the supporting evidence at January's recommended configuration,
where picking utilisation is identical across the three policies to three decimal places, so
capacity is held constant.

\begin{table}[htbp]
\centering
\caption[Picking waiting times at January \rgm{s13\_8\_4}]{Picking waiting times at January's
recommended configuration. Utilisation is identical across policies, so differences arise from
ordering alone.}
\label{tab:app-waits}
\small
\begin{tabular}{lrrrr}
\toprule
\textbf{Policy} & \textbf{Utilisation} & \textbf{Mean wait} & \textbf{p95 wait} &
\textbf{Normal SLA} \\
 & & \textbf{(min)} & \textbf{(min)} & \textbf{(\%)} \\
\midrule
FIFO         & 0.958 & 380.4 & 842.4  & 91.4 \\
Urgent-First & 0.958 & \textbf{301.8} & \textbf{1796.9} & 80.0 \\
RL-3         & 0.958 & 379.9 & 993.8  & 88.8 \\
\bottomrule
\end{tabular}
\end{table}

The pattern is distinctive. Urgent-First achieves the \emph{lowest} mean picking wait of the
three policies while producing by far the \emph{highest} p95 wait --- more than double FIFO's.
It serves the average order sooner while leaving a minority waiting far longer. Since the
service criterion is a threshold rather than a mean, the tail is what determines attainment,
and Urgent-First's normal-class attainment at this configuration (\SI{80.0}{\percent}) sits
exactly on its floor.

RL-3, whose sub-queues are first-in-first-out within each class, produces a mean wait almost
identical to FIFO's and a p95 between the two. This is consistent with the interpretation that
heap-based within-class ordering inflates the tail without changing the mean, and it is why
\Cref{sec:future-work} identifies a first-in-first-out Urgent-First variant as the most valuable
single extension to this study.

\section{Break-even analysis for the marginal FTE}

The economic test applied by the adaptive search can be expressed as a break-even quantity: how
many late orders must an additional worker prevent in order to pay for itself.

Under the retrospective run's assumptions, one FTE costs
\(\num{2400} = 160 \times 15\) euros per month. Against a late urgent order penalty of
\eur{20} and a late normal penalty of \eur{5}, break-even requires preventing:

\begin{center}
\begin{tabular}{lr}
\toprule
\textbf{If failures prevented were entirely\ldots} & \textbf{Orders required} \\
\midrule
urgent orders & 120 \\
normal orders & 480 \\
the observed December mix & approximately 456 \\
\bottomrule
\end{tabular}
\end{center}

The observed average penalty per late order in December was EUR~5.27, and the tested additional
picking worker prevented \eur{20} of penalties in total. The margin is therefore not close: the
additional worker would need to have been roughly two orders of magnitude more effective to
justify its cost.

\section{Reconciling the selection rules}
\label{sec:app-rules}

Two distinct things are computed per month and are easily conflated: the framework's
\emph{recommendation}, and a set of \emph{descriptive} summary statistics published alongside it
for comparison. \Cref{tab:app-rules} separates them.

\begin{table}[htbp]
\centering
\caption[Selection rules]{The recommendation rule and the descriptive statistics reported
alongside it.}
\label{tab:app-rules}
\small
\begin{tabularx}{\textwidth}{lX}
\toprule
\textbf{Rule} & \textbf{Definition} \\
\midrule
\textbf{Framework recommendation} & Lexicographic, as specified in \Cref{sec:selection-rule}:
  among configurations meeting \emph{both} service floors, the one of lowest total cost; if none
  meets both floors, the one of smallest service violation, ties broken by cost, reported as
  infeasible. In forecast mode only three-replication aggregates are eligible. \\
\midrule
\multicolumn{2}{l}{\emph{Descriptive statistics, reported for comparison --- not the
  recommendation rule}} \\
Lowest total cost & The configuration minimising total cost across all candidates and policies,
  with no feasibility filter. \\
Smallest feasible & Among configurations meeting both service floors, the one using the fewest
  total FTE. \\
Smallest meeting urgent floor & Among configurations meeting the urgent floor alone, the one
  using the fewest FTE. \\
Best per policy & For each policy separately, its own lowest-cost configuration. \\
\bottomrule
\end{tabularx}
\end{table}

\Cref{tab:recommendations} is tabulated from the lowest-total-cost column of the export. In this
run that column coincides with the framework's lexicographic recommendation in \emph{all twelve
months} --- configuration, policy and cost are identical row by row --- because a feasible
configuration existed in every month and, in every month, the globally cheapest configuration
happened also to be feasible. The two rules are nevertheless different rules, and a run
containing a cheap infeasible configuration would separate them: the descriptive column would
report it, and the recommendation would not.
